\documentclass[11pt]{article}

\usepackage[T1]{fontenc}
\usepackage{lmodern}
\usepackage[margin=1in]{geometry}
\usepackage{amsmath,amssymb,amsthm,mathtools}
\usepackage{enumitem}
\usepackage{booktabs,array}
\usepackage{xcolor}
\usepackage[protrusion=true,expansion=false]{microtype}
\usepackage{xurl}
\usepackage[colorlinks=true,linkcolor=blue,citecolor=blue,urlcolor=blue]{hyperref}

\hypersetup{
  pdftitle={A Formalized Proof of Cheeger--Gromov--Hamilton Compactness},
  pdfauthor={Bennett Chow, Yuan Liao, and Ziyang Qin}
}

\theoremstyle{plain}
\newtheorem{theorem}{Theorem}[section]
\newtheorem{lemma}[theorem]{Lemma}
\newtheorem{proposition}[theorem]{Proposition}
\newtheorem{corollary}[theorem]{Corollary}
\theoremstyle{definition}
\newtheorem{definition}[theorem]{Definition}
\theoremstyle{remark}
\newtheorem{remark}[theorem]{Remark}

\newcommand{\Rm}{\operatorname{Rm}}
\newcommand{\Ric}{\operatorname{Ric}}
\newcommand{\Vol}{\operatorname{Vol}}
\newcommand{\inj}{\operatorname{inj}}
\newcommand{\dist}{\operatorname{dist}}
\newcommand{\diam}{\operatorname{diam}}
\newcommand{\Id}{\operatorname{Id}}
\newcommand{\Exp}{\operatorname{Exp}}
\newcommand{\tr}{\operatorname{tr}}
\newcommand{\dd}{\,\mathrm d}
\newcommand{\eps}{\varepsilon}
\newcommand{\del}{\delta}
\newcommand{\N}{\mathbb N}
\newcommand{\R}{\mathbb R}
\newcommand{\norm}[1]{\left|#1\right|}
\newcommand{\abs}[1]{\left|#1\right|}
\newcommand{\ip}[2]{\left\langle #1,#2\right\rangle}
\newcommand{\leanref}[1]{\nolinkurl{#1}}
\newcommand{\leanfile}[1]{\nolinkurl{#1}}
\makeatletter
\g@addto@macro{\UrlBreaks}{\do\_\do\.\do\/}
\makeatother
\Urlmuskip=0mu plus 2mu\relax
\emergencystretch=3em

\title{A Formalized Proof of Cheeger--Gromov--Hamilton Compactness}
\author{Bennett Chow \and Yuan Liao \and Ziyang Qin}
\date{Working draft, July 2026}

\begin{document}
\maketitle

\begin{abstract}
We give a mathematical account of a Lean formalization of two compactness
theorems used in Ricci flow.  The first is pointed smooth
Cheeger--Gromov compactness for complete connected Riemannian manifolds with
uniformly bounded curvature derivatives and a uniform injectivity-radius
bound at the basepoints.  The second is Hamilton's compactness theorem for
complete pointed Ricci flows on a common open time interval, assuming
curvature bounds on compact time windows and a time-zero basepoint
injectivity-radius bound.

The proof is not recorded as a single appeal to a classical black box.
Starting from the geometric hypotheses, the formalization constructs
injectivity-radius decay, volume and packing estimates, controlled normal
coordinates, stabilized nets, approximate comparison diffeomorphisms, a
smooth direct-limit manifold, a complete limiting metric, and finally a
smooth limiting Ricci flow.  Special attention is paid to the points where a
paper proof silently identifies local branches, changes subsequences, or
passes from local convergence to completeness.  We also explain how the
completed compactness theorems enter the formalized Hamilton
positive-Ricci-curvature program and distinguish the compactness result from
the remaining uniform-existence frontier.
\end{abstract}

\noindent\textbf{Keywords.}
Ricci flow; Cheeger--Gromov compactness; Hamilton compactness; bounded
geometry; injectivity radius; smooth direct limits; formalized mathematics;
Lean.

\medskip
\noindent\textbf{Mathematics Subject Classification (2020).}
Primary 53E20, 53C20, 68V20; Secondary 35K59, 53C21.

\tableofcontents

\section{Introduction}

Compactness arguments in geometric analysis replace a sequence of spaces
whose geometry is uniformly controlled by a single limiting space on which
one can pass inequalities, differential equations, and curvature identities
to the limit.  In Ricci flow the relevant notion is smooth pointed
Cheeger--Gromov--Hamilton convergence.  One must extract both a limiting
manifold and comparison maps, prove smooth convergence of the pulled-back
metrics, and retain completeness of the limit.

In an informal proof, several operations are often suppressed:
\begin{itemize}[leftmargin=2em]
  \item one repeatedly passes to subsequences and silently replaces the old
  indices by the new ones;
  \item one chooses normal coordinates whose radius depends on the distance
  from a basepoint;
  \item one glues local maps which are only approximately compatible;
  \item one constructs a manifold from an increasing system of open sets;
  \item and one upgrades local smooth convergence to completeness of the
  limiting metric.
\end{itemize}
Each of these is mathematically standard, but each carries nontrivial domain,
compatibility, and quantifier information.  The purpose of the formalization
is to make that information explicit without changing the classical
argument.

The metric theorem corresponds to Theorem 3.9 of
Morgan--Tian \cite{MR2274812}.  The Ricci-flow theorem corresponds to their
Theorem 3.10 and to Hamilton's compactness theorem
\cite{HamiltonCompactness}.  The formal proof follows the bounded-geometry
route: construct controlled charts, compare distant members of the sequence
on larger and larger balls, form a direct limit, and then use Shi estimates
\cite{Shi1989a} to recover the space--time statement.

\subsection{Main results}

We first state the two results in ordinary mathematical language.

\begin{theorem}[Pointed smooth metric compactness]\label{thm:metric}
Let
\[
  (M_i,g_i,x_i),\qquad i\in\N,
\]
be complete connected pointed smooth Riemannian \(n\)-manifolds without
boundary.  Suppose that for every \(m\geq 0\) there is a constant
\(C_m<\infty\), independent of \(i\), such that
\[
  \sup_{M_i}\norm{\nabla^m\Rm(g_i)}_{g_i}\leq C_m.
\]
Suppose also that there is \(\iota_0>0\) such that
\[
  \inj_{g_i}(x_i)\geq \iota_0
  \qquad\text{for every }i.
\]
Then there are a subsequence, still denoted by \(i\), a complete connected
pointed smooth Riemannian manifold
\[
  (M_\infty,g_\infty,x_\infty),
\]
and pointed smooth embeddings from an exhaustion of \(M_\infty\) into the
members of the subsequence such that
\[
  (M_i,g_i,x_i)\longrightarrow
  (M_\infty,g_\infty,x_\infty)
\]
in the smooth pointed Cheeger--Gromov sense.
\end{theorem}

The formal conclusion uses open source domains rather than a preselected
compact exhaustion.  It produces increasing comparison domains whose union
is the limit and partial diffeomorphisms into the sequence members.  On every
compact subset \(K\subset M_\infty\) and for every \(p\geq0\),
\[
  \sup_K
  \norm{\nabla_{g_\infty}^{\,p}
    \bigl(\Phi_i^*g_i-g_\infty\bigr)}_{g_\infty}
  \longrightarrow 0.
\]
This formulation is equivalent to the usual compact-exhaustion version and
is better suited to keeping domains and restrictions explicit.

\begin{theorem}[Compactness for complete Ricci flows]\label{thm:flow}
Let \(J=(a,b)\subset\R\) be an open interval containing \(0\), and let
\[
  (M_i,g_i(t),x_i),\qquad t\in J,
\]
be complete connected pointed Ricci flows on \(n\)-manifolds without
boundary.  Assume that for every compact interval \(J'\Subset J\) there is
\(C(J')<\infty\) such that
\[
  \sup_i\sup_{M_i\times J'}\norm{\Rm(g_i(t))}_{g_i(t)}^2
  \leq C(J').
\]
Assume also that
\[
  \inj_{g_i(0)}(x_i)\geq\iota_0>0
\]
uniformly in \(i\).  Then a subsequence converges smoothly in the pointed
Cheeger--Gromov--Hamilton sense on \(J\) to a pointed Ricci flow
\[
  (M_\infty,g_\infty(t),x_\infty),\qquad t\in J.
\]
Every time slice \((M_\infty,g_\infty(t))\) is complete.
\end{theorem}

\begin{remark}[What is and is not assumed]
Theorem~\ref{thm:flow} assumes only a zeroth-order curvature bound on each
compact time window.  Higher spatial derivative estimates are produced from
the Ricci-flow equation by complete Shi estimates.  Thus the all-order
bounded geometry needed by Theorem~\ref{thm:metric} is a consequence at
time zero, not an additional hypothesis.
\end{remark}

\subsection{Formal status}

The formal counterparts of Theorems~\ref{thm:metric} and
\ref{thm:flow} are respectively \leanref{metricCompactness} and
\leanref{compactnessSol}.
In the July 2026 working tree both are checked without a project-defined
axiom or a \leanref{sorryAx} dependency.  Their direct axiom reports contain
only the standard logical principles \leanref{propext},
\leanref{Classical.choice}, and \leanref{Quot.sound}.  This status is
different from the earlier account in \leanfile{main24.tex}, where
Cheeger--Gromov--Hamilton compactness was still presented as a deferred
global input.

\section{Pointed convergence and bounded geometry}

\subsection{Pointed manifolds and convergence maps}

A pointed Riemannian manifold is a triple \((M,g,x)\).  A convergence datum
from a sequence \((M_i,g_i,x_i)\) to a pointed limit
\((M_\infty,g_\infty,x_\infty)\) consists of:
\begin{enumerate}[leftmargin=2em]
  \item a strictly increasing subsequence \(i_k\);
  \item increasing open sets \(U_k\subset M_\infty\) which contain
  \(x_\infty\) and exhaust \(M_\infty\);
  \item open sets \(V_k\subset M_{i_k}\);
  \item pointed diffeomorphisms
  \[
    \Phi_k:U_k\longrightarrow V_k,
    \qquad \Phi_k(x_\infty)=x_{i_k};
  \]
  \item compact-open \(C^\infty\) convergence
  \[
    \Phi_k^*g_{i_k}\longrightarrow g_\infty.
  \]
\end{enumerate}
The exhaustion property means that every compact set \(K\subset M_\infty\)
is contained in \(U_k\) for all sufficiently large \(k\).

The formalization packages the maps separately from the convergence of the
metrics.  This separation is important.  A partial diffeomorphism determines
where pullback makes sense, while metric convergence is a quantified
statement on compact subsets of its source.

\subsection{Uniform bounded geometry}

For a fixed pointed manifold, bounded geometry means that every covariant
derivative of the curvature tensor is globally bounded.  For a sequence we
require a single function
\[
  m\longmapsto C_m
\]
such that the \(m\)-th estimate holds for every member:
\[
  \norm{\nabla^m\Rm(g_i)}\leq C_m.
\]
No growth condition in \(m\) is imposed.  Compactness at a fixed derivative
order uses only finitely many of these constants; the final smooth
subsequence is obtained by a diagonal argument over all orders.

The tensor convention in the formal development uses the lowered
\((0,m+4)\)-tensor \(\nabla^m\Rm_{04}\).  This is equivalent to the usual
operator-valued formulation because the metric identifies the raised and
lowered versions.  It also makes the norm and pullback statements uniform in
the tensor valence.

\section{From basepoint control to controlled local geometry}

The central local-to-global issue is that a lower bound for
\(\inj_{g_i}(x_i)\) is given only at the distinguished point, while the
compactness construction needs charts at every point in larger and larger
balls.  The first part of the proof propagates a positive, though decaying,
injectivity scale away from the basepoint.

\subsection{Injectivity-radius decay}

\begin{proposition}[Controlled injectivity scale]\label{prop:injdecay}
Under the hypotheses of Theorem~\ref{thm:metric}, for every \(R<\infty\)
there is \(\lambda(R)>0\), depending only on \(R\), the dimension, the
curvature bounds, and \(\iota_0\), such that
\[
  y\in B_{g_i}(x_i,R)
  \quad\Longrightarrow\quad
  \inj_{g_i}(y)\geq\lambda(R)
\]
for every \(i\).
\end{proposition}

\begin{proof}[Proof outline]
Completeness and connectedness identify the intrinsic Riemannian distance
with the proper metric used to form balls and minimizing geodesics.  Curvature
control gives a uniform conjugate-radius estimate on bounded regions.  If
the injectivity radius degenerates at a point before the conjugate radius,
one obtains a short geodesic loop.

The Cheeger--Gromov--Taylor argument compares such loops with the volume of
controlled balls.  A localized Whitehead construction supplies unique short
inverse branches of the exponential map.  Strict convexity of the half
squared-distance function, proved through the Jacobi-field Hessian formula,
prevents incompatible short branches.  The resulting propeller-count
estimate bounds the number of short lifts.  Combined with the lower volume
available at the basepoint and relative volume comparison along a chain of
balls, this rules out arbitrarily small injectivity radius on
\(B_{g_i}(x_i,R)\).

The argument produces an explicit positive profile \(\lambda(R)\).  The
profile may decay as \(R\to\infty\); only positivity for each fixed \(R\) is
needed.
\end{proof}

The proof must use the actual Riemannian distance.  Replacing it by an
arbitrary stored distance would make the volume and ball statements
unusable.  In the formalization the realization theorem identifies the
extended distance in the decay package with the metric distance of each
sequence member.

\subsection{Volume comparison and packing}

The curvature bound gives a Ricci lower bound, so Bishop--Gromov comparison
controls the relative volumes of balls.  Together with
Proposition~\ref{prop:injdecay}, it yields both a lower bound for small balls
and an upper bound for the number of separated points in a bounded ball.

\begin{lemma}[Finite packing on bounded balls]\label{lem:packing}
Fix \(R<\infty\) and \(0<r\leq\lambda(R)\).  There is a finite number
\(A(R,r)\), independent of \(i\), such that every \(r\)-separated subset of
\(B_{g_i}(x_i,R)\) has cardinality at most \(A(R,r)\).
\end{lemma}

\begin{proof}
The balls of radius \(r/2\) centered at points of an \(r\)-separated family
are pairwise disjoint.  Injectivity and curvature control give a uniform
lower bound for the volume of each small ball.  Bishop--Gromov comparison
places all of them in a slightly larger ball with uniformly controlled
volume.  Comparing the total volume of the disjoint union with the containing
ball gives the cardinality estimate.
\end{proof}

In the implemented route the containing scale is capped explicitly.  This
matters because a statement uniform in an unbounded ratio of radii would be
false, for example on hyperbolic space.  The cap is threaded through every
consumer which uses the packing estimate.

\subsection{Stable nets}

For each large ball choose a maximal separated set.  Maximality makes the
corresponding enlarged balls a cover, while Lemma~\ref{lem:packing} bounds the
number of centers and the overlap multiplicity.  Since only finitely many
incidence patterns can occur at each fixed radius, repeated finite
subsequence extraction stabilizes:
\begin{itemize}[leftmargin=2em]
  \item the number of centers;
  \item which controlled coordinate balls overlap;
  \item the parent relation between consecutive radii;
  \item and the finite combinatorial data needed to define transition maps.
\end{itemize}
A diagonal subsequence makes these choices compatible for every integer
radius.

\section{Uniform normal coordinates}

\subsection{A branch-carrying normal-coordinate package}

At a center \(y\in B_{g_i}(x_i,R)\), choose a radius
\[
  0<\rho_R<
  \min\{\lambda(R),\rho_{\mathrm{Jac}}(R)\},
\]
where \(\rho_{\mathrm{Jac}}(R)\) is a uniform radius on which Jacobi-field
estimates control the differential of the exponential map.  The exponential
map then has a genuine inverse branch on the whole controlled ball.

\begin{proposition}[Uniform normal-coordinate control]\label{prop:h6}
For every finite \(R\) there is a radius \(\rho_R>0\) such that each center
\(y\in B_{g_i}(x_i,R)\) admits a normal chart
\[
  \chi_{i,y}:B(0,\rho_R)\subset T_yM_i\longrightarrow M_i
\]
with the following properties.
\begin{enumerate}[leftmargin=2em]
  \item The chart is the selected inverse branch of the exponential map on
  its image.
  \item In Euclidean coordinates,
  \[
    \frac12\del_{ab}\leq (\chi_{i,y}^*g_i)_{ab}\leq2\del_{ab}
  \]
  as quadratic forms.
  \item For every \(m\geq0\) there is \(C'_m(R)<\infty\), independent of
  \(i\) and \(y\), such that
  \[
    \sup_{B(0,\rho_R)}
    \norm{\partial^m(\chi_{i,y}^*g_i)}
    \leq C'_m(R).
  \]
  \item On overlaps, the transition maps and their inverses have uniform
  derivative bounds of every prescribed finite order.
\end{enumerate}
\end{proposition}

\begin{proof}[Proof outline]
The differential of \(\Exp_y\) is described by Jacobi fields.  Curvature and
curvature-derivative bounds give uniform ODE estimates for the Jacobi
equation and its derivatives with respect to the initial vector.  On a
sufficiently small ball, \(D\Exp_y\) is close to the identity, hence
invertible with controlled inverse.  The injectivity-radius bound ensures
that the local inverse agrees with the selected intrinsic branch on the
whole ball.  Pulling back the metric and differentiating the Jacobi equation
gives the metric estimates.  Transition maps are compositions of one branch
with the inverse of another, so the same estimates and the inverse function
theorem control their derivatives.
\end{proof}

The branch is part of the data.  Merely knowing that some chart exists at
each point is insufficient: later overlap and cocycle statements must refer
to the same chart.  This is one of the main proof-engineering differences
between the formal argument and a short textbook proof.

\section{Approximate comparison maps}

\subsection{Local limits of chart metrics and transitions}

Fix a controlled radius level.  There are finitely many net centers and
finitely many relevant overlaps.  Proposition~\ref{prop:h6} gives uniform
\(C^m\) bounds on all chart metrics and transition maps.  Arzela--Ascoli and
a diagonal argument therefore give smooth limits on slightly smaller balls.
The limiting transition maps retain the inverse and cocycle identities
because these identities hold before taking the limit and composition is
continuous in the relevant compact-open topology.

The finite chart system defines a smooth model for the bounded region.  The
remaining task is to compare two sufficiently late sequence members on that
region by a single smooth map.

\subsection{Gluing by a center-of-mass construction}

Local chart identifications agree only approximately.  Let
\[
  f_\alpha:U_\alpha\subset M_i\longrightarrow M_j
\]
denote the local comparison maps obtained from corresponding chart centers,
and let \(\{\eta_\alpha\}\) be a partition of unity subordinate to the source
cover.  For \(x\) in the controlled source region, define a point
\(\Phi_{ij}(x)\) as the unique minimizer of
\[
  F_x(z)
  =
  \frac12\sum_\alpha
  \eta_\alpha(x)\,
  \dist_{g_j}\bigl(z,f_\alpha(x)\bigr)^2.
\]
All active points \(f_\alpha(x)\) lie in a common strongly convex ball.  The
Hessian of \(F_x\) is therefore positive definite.  The implicit function
theorem gives a smooth center \(\Phi_{ij}(x)\).

The uniform overlap bound from Lemma~\ref{lem:packing} prevents the sums from
growing with \(i\) or \(j\).  Differentiating the center equation and using
the uniform inverse Hessian estimate controls all derivatives of
\(\Phi_{ij}\).  Repeating the construction in the reverse direction and
comparing both compositions with the identity yields a genuine partial
diffeomorphism on a slightly smaller source ball.

\begin{proposition}[Textbook comparison step]\label{prop:b1}
After passing to one subsequence, the following holds.  Given \(r>0\),
\(0<\eps<1\), and \(p\in\N\), there is \(k_0\) such that for every
\(k,\ell\geq k_0\) there is a pointed partial diffeomorphism
\[
  \Phi_{k\ell}:M_k\dashrightarrow M_\ell
\]
whose source contains \(\overline{B}_{g_k}(x_k,r)\) and such that
\(\Phi_{k\ell}\) is an \(\eps\)-approximate isometry through order \(p\) on
that ball:
\[
  \sup_{\overline{B}_{g_k}(x_k,r)}
  \norm{\nabla_{g_k}^{\,q}
  \bigl(\Phi_{k\ell}^*g_\ell-g_k\bigr)}_{g_k}
  <\eps,
  \qquad 0\leq q\leq p.
\]
The inverse branch satisfies the analogous estimate on the image.
\end{proposition}

\begin{proof}
Choose a chart radius small compared with the injectivity and Jacobi radii on
a ball slightly larger than \(r\).  Stabilize the finite covering and overlap
data, take smooth limits of all chart metrics and transition maps, and form
the center-of-mass map described above.  The positive Hessian and the
implicit derivative estimates give the required \(C^p\) bounds.  The
slightly larger ambient radius absorbs the support of the partition of unity
and guarantees that the inverse construction is defined wherever it is
used.  Finally choose the subsequence far enough out that every error is
below \(\eps\).
\end{proof}

The formal theorem corresponding to Proposition~\ref{prop:b1} is
\leanref{MetricCompactBase.exists_b1}.  The implementation first proves a
stronger raw package which retains the forward and reverse comparison maps,
their domains, and their common chart provider; the textbook statement is a
projection of that package.

\section{The smooth direct limit}

\subsection{A coherent tail system}

Apply Proposition~\ref{prop:b1} with radii tending to infinity, errors tending
to zero, and derivative orders tending to infinity.  A further diagonal
subsequence gives adjacent comparison maps
\[
  \Psi_j:U_j\subset M_j\longrightarrow M_{j+1}
\]
defined on balls whose radii grow at least exponentially.  Restricting to
smaller balls \(B_j\Subset U_j\) gives:
\[
  \Psi_j(B_j)\Subset B_{j+1}.
\]
The restrictions form a smooth sequential system.  Its direct limit
\[
  M_\infty=\varinjlim(B_j,\Psi_j)
\]
is a smooth Hausdorff manifold.  All images of the basepoints represent one
point \(x_\infty\).

The use of smaller tail balls is essential.  The large balls are needed to
define the comparison maps; the gap between the large and small radii
provides relatively compact image inclusions and room for later path and
metric estimates.

\subsection{Limit metrics on the stages}

Fix a source stage \(B_j\).  Pull every later metric back to \(B_j\) by the
finite chain
\[
  \Psi_{j,\ell}
  =
  \Psi_{\ell-1}\circ\cdots\circ\Psi_j.
\]
The approximate-isometry estimates imply that
\[
  \Psi_{j,\ell}^*g_\ell
\]
is Cauchy in every \(C^p\) seminorm on compact subsets of \(B_j\).  After a
common diagonal extraction it converges smoothly to a metric
\(g_{\infty,j}\).

The chain identity gives the cocycle
\[
  \Psi_j^*g_{\infty,j+1}=g_{\infty,j}.
\]
Hence the stage metrics glue to a smooth Riemannian metric \(g_\infty\) on
the direct limit.  The order-zero estimates ensure that the limit is positive
definite; positivity is not inferred merely from pointwise convergence of
uncontrolled bilinear forms.

\subsection{Completeness}

Local smooth convergence does not by itself imply completeness.  The proof
uses compact cores inside the tail stages.  Let
\[
  K_j
  =
  \overline{B}_{g_j}(x_j,R_j/2)
  \subset B_{g_j}(x_j,R_j)=B_j.
\]
Completeness and connectedness of \(M_j\) give properness, so \(K_j\) is
compact.  Its image in \(M_\infty\) is compact.

The order-zero comparison estimate gives a uniform lower bound
\[
  g_\infty\geq \frac12 g_j
\]
on the relevant stage.  Consequently a path starting near the common
basepoint and leaving the image of \(K_j\) has \(g_\infty\)-length bounded
below by a fixed multiple of \(R_j\).  Since \(R_j\to\infty\), every finite
\(g_\infty\)-ball about \(x_\infty\) lies in one compact core.  Thus closed
bounded balls are compact, and the induced metric is complete.

\begin{lemma}[Exhaustion by compact cores]\label{lem:core}
For every \(R<\infty\) there is \(j\) such that
\[
  \overline{B}_{g_\infty}(x_\infty,R)
\]
is contained in the image of the compact core \(K_j\).
\end{lemma}

\begin{corollary}
The metric \(g_\infty\) is complete.
\end{corollary}

\begin{proof}
Lemma~\ref{lem:core} makes every Cauchy sequence eventually lie in one
compact set.  It therefore has a convergent subsequence, and the Cauchy
property forces the whole sequence to converge.
\end{proof}

\subsection{Convergence to the original sequence}

The direct-limit inclusions provide maps from each stage into \(M_\infty\).
Inverting them on their images and composing with the finite comparison
chains gives pointed maps
\[
  \Phi_j:U_j^\infty\longrightarrow M_{i_j}
\]
whose targets are the original ambient sequence members, not merely the open
tail stages.  The stagewise metric limits imply
\[
  \Phi_j^*g_{i_j}\longrightarrow g_\infty
\]
smoothly on compact subsets.

Every \(B_j\) is path connected, and the transition maps preserve the
basepoint.  The direct limit is therefore connected.  This proves
Theorem~\ref{thm:metric}.

\section{From metric compactness to flow compactness}

\subsection{Complete Shi estimates}

Let \(J'\Subset J\).  Choose a slightly larger compact interval
\[
  J'\Subset J''\Subset J.
\]
The curvature bound on \(J''\), completeness of every time slice, and the
Ricci-flow equation give, for every \(m\geq1\),
\[
  \sup_i\sup_{M_i\times J'}
  \norm{\nabla^m\Rm(g_i(t))}
  \leq C_m(J',J'').
\]
The time buffer is essential: derivative estimates at a time \(t\) consume
curvature control on a larger parabolic neighborhood.

The formal proof uses an evolving distance barrier and a point-centered
cutoff.  The cutoff is inserted in a Bernstein argument for the curvature
tower.  Kato-type inequalities and the Ricci-flow evolution of
\(\nabla^m\Rm\) control the error at the moving boundary.  This gives a
constants-first complete Shi estimate and then a uniform version for every
canonical compact subinterval of \(J\).

\subsection{The time-zero metric limit}

At \(t=0\), the Shi bounds yield
\[
  \sup_i\sup_{M_i}\norm{\nabla^m\Rm(g_i(0))}\leq C_m
\]
for every \(m\).  Theorem~\ref{thm:metric} therefore gives a subsequence, a
complete pointed limit \((M_\infty,g_\infty(0),x_\infty)\), and pointed
comparison maps \(\Phi_i\).

The maps are chosen once at time zero.  They are then used to pull back the
entire family:
\[
  \widetilde g_i(t)=\Phi_i^*g_i(t).
\]
This fixed-map formulation avoids differentiating a moving family of
comparison maps.

\subsection{Space--time compactness}

On each compact \(K\subset M_\infty\) and compact interval \(J'\Subset J\),
the spatial curvature-derivative bounds control the spatial derivatives of
\(\widetilde g_i(t)\).  The Ricci-flow equation
\[
  \partial_t\widetilde g_i(t)
  =
  -2\Ric(\widetilde g_i(t))
\]
converts spatial bounds into time-derivative bounds.  Repeatedly commuting
\(\partial_t\) with covariant derivatives yields uniform estimates for every
mixed derivative on \(K\times J'\).

Arzela--Ascoli and a diagonal argument over:
\[
  K_1\Subset K_2\Subset\cdots,\qquad
  J_1\Subset J_2\Subset\cdots\Subset J,
\]
and all derivative orders produce a smooth family \(g_\infty(t)\) such that
\[
  \Phi_i^*g_i(t)\longrightarrow g_\infty(t)
\]
in \(C^\infty_{\mathrm{loc}}(M_\infty\times J)\).
Passing to the limit in the Ricci-flow equation gives
\[
  \partial_tg_\infty=-2\Ric(g_\infty).
\]

\subsection{Completeness of all time slices}

Completeness at \(t=0\) comes from the metric compactness theorem.  On a
compact time interval \(J'\Subset J\), the curvature bound controls the
Ricci tensor and hence the time variation of vector lengths:
\[
  \left|\frac{\dd}{\dd t}\log g_\infty(t)(v,v)\right|
  \leq C(J').
\]
Integrating gives uniform metric equivalence
\[
  e^{-C|t-s|}g_\infty(s)
  \leq g_\infty(t)
  \leq e^{C|t-s|}g_\infty(s).
\]
Metrics which are uniformly equivalent have the same Cauchy sequences.
Therefore completeness propagates from \(t=0\) to every \(t\in J\).  This
finishes the proof of Theorem~\ref{thm:flow}.

\section{Correspondence with the formal proof}

The preceding proof is organized mathematically.  For reproducibility, this
section records the principal formal counterparts without turning the paper
into a line-by-line Lean commentary.

\begin{center}
\begin{tabular}{>{\raggedright\arraybackslash}p{0.31\textwidth}
                >{\raggedright\arraybackslash}p{0.29\textwidth}
                >{\raggedright\arraybackslash}p{0.31\textwidth}}
\toprule
Mathematical ingredient & Principal declaration & Source module \\
\midrule
Uniform curvature derivatives &
\leanref{SeqBoundedGeometry} &
\leanfile{BoundedGeometry.lean} \\
Basepoint injectivity bound &
\leanref{BaseInjBound} &
\leanfile{InjectivityRadius.lean} \\
Injectivity decay &
\leanref{injDecay_of_bg} &
\leanfile{C4/CGTDecay.lean} \\
Volume and packing &
\leanref{volInput_of_bg}, \leanref{packInput_of_bg} &
\leanfile{C4/VolumeOverlap.lean} \\
Controlled normal charts &
\leanref{exists_h6NormalData} &
\leanfile{C4/H6NormalData.lean} \\
Pairwise approximate maps &
\leanref{MetricCompactBase.exists_b1} &
\leanfile{C4/StepB1Textbook.lean} \\
Direct limit and complete metric &
\leanref{compactness_canon} &
\leanfile{C4/StepDAssembly.lean} \\
Metric compactness &
\leanref{metricCompactness} &
\leanfile{C4/MetricCompactnessUncondH6.lean} \\
Complete Shi estimates &
\leanref{movingShi_open} &
\leanfile{MovingShiOpen.lean} \\
Ricci-flow compactness &
\leanref{compactnessSol} &
\leanfile{HamiltonCompactness.lean} \\
\bottomrule
\end{tabular}
\end{center}

\subsection{Why the conclusion is structured}

The metric theorem returns more than an existential statement.  Its
conclusion contains:
\begin{itemize}[leftmargin=2em]
  \item the subsequence and its strict monotonicity;
  \item the pointed limit manifold and completeness of its metric;
  \item the source and target domains of every comparison map;
  \item the partial diffeomorphisms themselves;
  \item and compact-open convergence at every derivative order.
\end{itemize}
This structure is what later curvature-transfer and Ricci-flow arguments
consume.  An opaque statement that merely says ``a convergent subsequence
exists'' would be too weak for downstream formalization.

\subsection{Subsequence bookkeeping}

The construction uses several diagonals: one for finite packing data, one for
chart transitions, one for derivative orders, one for direct-limit stage
metrics, and one for compact time windows.  The formal proof represents each
subsequence by an explicit strictly monotone function \(\N\to\N\) and proves
that compositions remain strictly monotone.  At the endpoint the convergence
conclusion is transported back through the nested subsequence maps to the
original sequence.

\subsection{Domains and actual branches}

Every coordinate comparison is restricted to its genuine open source.
Normal coordinates retain the selected inverse branch of the exponential
map, and every overlap theorem carries the membership hypotheses needed to
use it.  This prevents two common informal ambiguities:
\begin{enumerate}[leftmargin=2em]
  \item evaluating a local frame or inverse branch outside its domain;
  \item using two existentially chosen charts as if they were definitionally
  the same chart.
\end{enumerate}

\section{Role in the Hamilton positive-Ricci program}

The compactness theorem is used after selecting and rescaling a sequence of
high-curvature points.  The rescaled flows share a fixed time window, have
uniform curvature bounds there, and have a uniform time-zero injectivity
bound obtained from the proved noncollapsing and volume-to-injectivity
machinery.  Theorem~\ref{thm:flow} then gives a complete smooth pointed limit.

The provider-native Hamilton adapter uses this limit directly.  It no longer
depends on the legacy theorem-shaped placeholder
\leanref{ham3_cgh_limit}.  Curvature and pinching transfer are derived from
the actual smooth convergence witness, the limit is shown to have constant
positive sectional curvature, and compactness plus the global convergence
maps transfer the resulting metric back to the original manifold.

\begin{remark}[Separation of compactness from the Hamilton endpoint]
Closing Theorems~\ref{thm:metric} and \ref{thm:flow} does not by itself close
the entire Hamilton theorem.  Compactness begins only after a sufficiently
long Ricci flow and a normalized blow-up sequence have been constructed.
The current provider-native Hamilton endpoint has one remaining nonstandard
axiom source: the uniform-existence producer used to construct the required
maximal flow package.
\end{remark}

\section{Remaining formalization boundary}

\subsection{No deferred compactness input on the new route}

The metric compactness theorem, the complete Shi estimates, the Ricci-flow
compactness theorem, the volume-to-injectivity bridge, and the Hamilton
compactness adapter are all proved on the current provider-native route.
Accordingly, the final presentation should not list
Cheeger--Gromov--Hamilton compactness, Shi estimates, Perelman
noncollapsing, or the spherical-space-form classification as deferred inputs
of that route.

\subsection{The live critical frontier}

The remaining critical input is the uniform Ricci-flow existence theorem
represented by \leanref{ham3_flow_exists_normalized}.
Its role is to construct a common positive existence time from uniform
ellipticity and finitely many initial metric-jet bounds, and then to support
the continuation argument.  The analytic work is concentrated in a
time-dependent Ricci--DeTurck operator:
\begin{itemize}[leftmargin=2em]
  \item a small second-order perturbation acting on the \(H^4\to H^2\)
  scale;
  \item a first-order term acting on the adjacent
  \(H^3\to H^2\) and \(H^2\to H^1\) scales;
  \item strong measurability and time-integrability of these operator
  families along a rough solution;
  \item a same-horizon nonautonomous maximal-regularity bootstrap;
  \item and all-order smoothing, including compatibility at \(t=0\).
\end{itemize}
The static low-order operator decomposition and its Sobolev bounds are
substantially developed, but the time-dependent packaging and final
same-horizon assembly remain the active frontier.

\subsection{Noncritical legacy and library frontiers}

Several unfinished declarations are not dependencies of the completed
compactness theorem or of the new Hamilton adapter:
\begin{itemize}[leftmargin=2em]
  \item the legacy \leanref{ham3_cgh_limit} wrapper still has an unfinished
  proof, but the provider-native Hamilton route bypasses it;
  \item a generic arbitrary-valence Weyl eigenvalue-counting theorem remains
  deferred, while the live scalar spectral consumers use a proved
  Weyl-free tail argument;
  \item a textbook-style Hopf--Rinow capstone file remains incomplete, while
  the properness and minimizing-geodesic results used by compactness are
  proved through a separate route.
\end{itemize}
These are legitimate cleanup or library-completion projects, but they should
not be displayed as blockers for Theorems~\ref{thm:metric} or
\ref{thm:flow}.

\section{Recommended final presentation}

For a final paper or research talk, the compactness result is clearest when
presented as a theorem with three mathematical innovations in the
formalization, rather than as a chronology of file construction.

\paragraph{First: state the two closed theorems.}
Begin with Theorems~\ref{thm:metric} and \ref{thm:flow}, including
completeness, connectedness, the common open time interval, and the exact
curvature hypotheses.

\paragraph{Second: isolate the three proof bottlenecks.}
The mathematical body of the presentation should emphasize:
\begin{enumerate}[leftmargin=2em]
  \item propagation from basepoint injectivity to a controlled local scale,
  using CGT, volume comparison, and packing;
  \item construction of one coherent chart provider and pairwise
  approximate diffeomorphisms by center-of-mass gluing;
  \item the smooth direct limit and the compact-core proof of completeness.
\end{enumerate}
The Ricci-flow upgrade can then be presented as a fourth, shorter step:
complete Shi estimates plus a space--time diagonal.

\paragraph{Third: explain formalization-specific choices after the proof.}
Only after the mathematical proof should one discuss explicit subsequences,
actual branch domains, dependent tensor fibers, and structured convergence
data.  These choices explain why the Lean development is large without
obscuring the geometry.

\paragraph{Fourth: end with an honest dependency ledger.}
The ledger should show metric compactness, flow compactness,
noncollapsing, and the Hamilton compactness adapter as completed.  The
uniform-existence theorem should be the sole red node on the new
Hamilton-main dependency path.  Legacy unfinished wrappers and independent
library generalizations should appear in a separate gray ``not on the
critical path'' box.

\section{Conclusion}

The formal proof follows the classical compactness strategy but makes its
hidden coherence data explicit.  Uniform curvature control and one
basepoint injectivity bound produce local geometric control on every bounded
region.  Stable nets and controlled normal charts yield smooth comparison
maps between late sequence members.  These maps define a smooth direct
limit, compact cores prove completeness, and the resulting comparison maps
give pointed \(C^\infty\) convergence.  Complete Shi estimates and the
Ricci-flow equation upgrade the static theorem to smooth
Cheeger--Gromov--Hamilton compactness on a common open time interval.

The compactness portion of the Hamilton program is therefore no longer a
deferred input on the provider-native route.  The remaining critical work is
analytic rather than compactness-theoretic: a uniform short-time and
continuation package for the Ricci--DeTurck construction.  Presenting the
project in this order makes the mathematical result visible while preserving
an exact account of what the final Hamilton endpoint still assumes.

\bibliographystyle{amsalpha}
\bibliography{references}

\end{document}
